\section{子任务 5：第 4 章 Decidability}

\subsection{核心知识点}

\begin{itemize}[leftmargin=2em]
  \item 可判定语言：存在总会停机并给出正确是非答案的图灵机。
  \item 对正则语言与上下文无关语言的若干判定问题，可以构造显式算法解决。
  \item 对角化方法：通过“与第 $i$ 台机器在第 $i$ 个输入上作对”制造不可判定性。
  \item 存在不可判定语言，也存在甚至不可识别的语言。
\end{itemize}

\subsection{典型问题与证明套路}

\begin{enumerate}[leftmargin=2em]
  \item \textbf{证明某问题可判定}：寻找有限搜索空间。
  
  对 DFA，常把问题转成图搜索或积自动机；对 CFG，常先规整化，再检查可生成性、可达性或有限长度推导。
  \item \textbf{证明某问题不可判定}：构造自指机器或借助对角化语言，证明任何声称“全都能判”的机器都会在某个特别输入上出错。
  \item \textbf{证明某语言不可识别}：往往把识别器与补语言的识别器合在一起，推出本应不可判定的问题却被判定，从而矛盾。
\end{enumerate}

\subsection{证明背后的哲学}

这一章第一次把“能力边界”讲得非常锋利：\textbf{不是我们暂时没想到算法，而是根本不存在算法。}
这里的不可能性来源于语义层面的自指。只要机器能够接收“机器描述”作为输入，就有机会构造“这台机器在谈论它自己”的局面，而对角化恰恰利用这种自指击穿任何全能判定器。

\subsection{本章精华}

\begin{itemize}[leftmargin=2em]
  \item 可判定性的证明常依赖“有限性”；不可判定性的证明常依赖“自指性”。
  \item 从这一章开始，理论计算不再只是分类语言，而是开始讨论\textbf{知识与算法的极限}。
\end{itemize}

\newpage
